%!TEX root = ../../main.tex
\section{정의 민감도: 계보의 증거와 미실행 항목}
\label{sec:eng-sensitivity}

M-RQ1의 두 번째 절은 경계와 규모 기준의 선택이 사례 구성에 어떤 변화를 만드는가이다. 이 장이 보고한 증거는 두 층으로 나뉜다. 첫째, 사례 집합 자체의 변화이다. $G$의 plateau(ARI $\ge 0.9$), $D$의 좁은 구간과 지역별 편차 7\% 이내, 게이트에 따른 인정 비율의 여섯 배 변화, 절단에 따른 규모 계급의 부호 반전은 모두 현재 코퍼스 또는 그 파일럿에서 측정되었다. 둘째, 예측 결과의 변화이다. $G \times D$ 스윕(공통 교전에서 $-0.0094$에서 $+0.0013$), 게이트 변경($+0.0039$, $-0.0066$), 관측 창 길이 스윕은 선행 연구 계보의 교전 승자 라벨과 LightGBM 학습기 아래에서 측정되었고, 그 계보에서 정의의 예측 타당성은 \numEngLineage{} 교전에 대한 경기 단위 out-of-fold AUC \aucLineageOOF{}로 요약된다.

두 번째 층은 본 논문의 주 목표($\SVI$, 동결 $\Vhat$, $q$ 대 $\PTflex$)에서 재실행되지 않았다. 보조 상수(앵커 위치, 상점 사건 제외, 라벨 지평, 처형 킬, 위치 격자 오차)의 민감도도 선행 연구 시기나 파일럿 코퍼스에서 측정된 것이다. 이 재실행은 부록 \ref{app:deferred}의 F5(정의 상수 민감도)로 이연되며, 본 논문은 이를 ``보였다''고 쓰지 않는다. 사람의 독립적 경계 주석이 없으므로 실제 한타 탐지 정확도도 주장하지 않는다.
